
\documentclass[addpoints]{exam}
\usepackage{amsmath}
\usepackage{amssymb}
\usepackage{color}
\usepackage{booktabs}

% \printanswers

\newcommand{\event}{Probability}
\newcommand{\tournament}{}
\def\type{0}

\setlength\fillinlinelength{\textwidth}
\CorrectChoiceEmphasis{\color{red}\bfseries}
\SolutionEmphasis{\color{red}\bfseries}

\setlength\answerclearance{2pt}
\pagestyle{headandfoot}
\runningheadrule
\runningheader{\event}{\tournament}{\ifnum\type=1 Class Set Test \else Full Name:\hspace{1cm} \fi}
\runningfootrule
\runningfooter{}{Page \thepage\ of \numpages}{}

\begin{document}
\begin{center}
    {\huge Grade 12 Data Management \\ \event
    \ifprintanswers \space Key
    \else \space \ifnum\type<2 Test \fi \ifnum\type=1 \textbf{(CLASS SET)} \fi
          \ifnum\type=2 Answer Sheet \fi
    \fi \\ \vspace{3mm}\tournament\vspace{1mm}}
\end{center}

\vspace{.25\textheight}

\begin{center}
    First Name: \ifprintanswers
    \underline{\hspace{3.5cm}}\textbf{KEY}\underline{\hspace{5.5cm}}\\\vspace{5mm}
    \else \underline{\hspace{10cm}}\\ \vspace{5mm} \fi
    Last Name: \underline{\hspace{10cm}}\\ \vspace{5mm}
\end{center}

\noindent\textbf{\underline{Directions:}}
\begin{itemize}
    \ifnum\type=1 \item \textbf{This test is a class set.} Do not write on this paper. \fi
    \item Show all work for full marks. Express probabilities as fractions or decimals (3 d.p.).
    \item The test is designed to be completed in 75 minutes.
\end{itemize}
\begin{center}
\textit{For grading use only}\\
\multirowgradetable{1}[pages]
\end{center}

\newpage
\begin{questions}

\section*{Part A --- Multiple Choice (15 marks)}
\vspace{5mm}

\question[1] If $P(A) = 0.4$, then $P(A^c)$ equals:
\begin{choices}
    \choice $0.4$
    \choice $0.16$
    \correctchoice $0.6$
    \choice $1.4$
\end{choices}

\question[1] Two events are mutually exclusive if:
\begin{choices}
    \choice $P(A \cap B) = P(A) \cdot P(B)$
    \correctchoice $P(A \cap B) = 0$
    \choice $P(A \cup B) = 1$
    \choice $P(A) = P(B)$
\end{choices}

\question[1] For any two events $A$ and $B$, the addition rule states:
\begin{choices}
    \choice $P(A \cup B) = P(A) + P(B)$
    \choice $P(A \cup B) = P(A) \cdot P(B)$
    \correctchoice $P(A \cup B) = P(A) + P(B) - P(A \cap B)$
    \choice $P(A \cup B) = P(A) - P(B)$
\end{choices}

\question[1] A card is drawn from a standard 52-card deck. The probability of drawing a
heart or a face card is:
\begin{choices}
    \choice $\dfrac{16}{52}$
    \choice $\dfrac{25}{52}$
    \correctchoice $\dfrac{22}{52}$
    \choice $\dfrac{3}{52}$
\end{choices}

\question[1] Events $A$ and $B$ are independent if and only if:
\begin{choices}
    \choice $P(A \cup B) = 1$
    \choice $P(A \cap B) = 0$
    \correctchoice $P(A \cap B) = P(A) \cdot P(B)$
    \choice $P(A) = P(B)$
\end{choices}

\question[1] The conditional probability $P(A \mid B)$ is defined as:
\begin{choices}
    \correctchoice $\dfrac{P(A \cap B)}{P(B)}$
    \choice $\dfrac{P(A)}{P(B)}$
    \choice $P(A) \cdot P(B)$
    \choice $P(A) + P(B) - 1$
\end{choices}

\question[1] A bag contains 3 red and 5 blue marbles. Two are drawn without replacement.
The probability both are red is:
\begin{choices}
    \choice $\dfrac{9}{64}$
    \correctchoice $\dfrac{3}{28}$
    \choice $\dfrac{6}{56}$
    \choice $\dfrac{1}{8}$
\end{choices}

\question[1] A fair die is rolled twice. The probability of getting a 6 on both rolls is:
\begin{choices}
    \choice $\dfrac{1}{6}$
    \choice $\dfrac{2}{6}$
    \correctchoice $\dfrac{1}{36}$
    \choice $\dfrac{11}{36}$
\end{choices}

\question[1] In a class, 60\% of students passed math, 70\% passed science, and 45\% passed
both. The probability a student passed math \textbf{given} they passed science is:
\begin{choices}
    \choice $0.45$
    \correctchoice $\dfrac{45}{70} \approx 0.643$
    \choice $\dfrac{45}{60} = 0.75$
    \choice $0.85$
\end{choices}

\question[1] The odds \textbf{in favour} of an event with probability $\dfrac{3}{8}$ are:
\begin{choices}
    \choice $3:8$
    \correctchoice $3:5$
    \choice $5:3$
    \choice $8:3$
\end{choices}

\question[1] A fair coin is flipped 3 times. The probability of getting \textbf{at least one head} is:
\begin{choices}
    \choice $\dfrac{1}{2}$
    \choice $\dfrac{3}{8}$
    \choice $\dfrac{4}{8}$
    \correctchoice $\dfrac{7}{8}$
\end{choices}

\question[1] $P(A \mid B) = P(A)$ is the condition for:
\begin{choices}
    \choice Mutual exclusivity
    \correctchoice Independence of $A$ and $B$
    \choice $A$ and $B$ being complementary
    \choice $P(B) = 0$
\end{choices}

\question[1] A committee of 2 is chosen at random from 4 men and 3 women. The probability
the committee has exactly 1 man is:
\begin{choices}
    \choice $\dfrac{4}{7}$
    \correctchoice $\dfrac{12}{21}$
    \choice $\dfrac{6}{21}$
    \choice $\dfrac{3}{7}$
\end{choices}

\question[1] A point is chosen at random inside a square of side 4. A circle of radius 1 is
inscribed at the centre. The probability the point lands inside the circle is:
\begin{choices}
    \choice $\dfrac{1}{4}$
    \correctchoice $\dfrac{\pi}{16}$
    \choice $\dfrac{\pi}{4}$
    \choice $\dfrac{1}{16}$
\end{choices}

\question[1] Suppose $P(B \mid A) = 0.8$ and $P(A) = 0.3$. Then $P(A \cap B) =$
\begin{choices}
    \choice $0.8$
    \choice $0.3$
    \correctchoice $0.24$
    \choice $0.38$
\end{choices}


\section*{Part B --- Short Answer (33 marks)}

\question A standard 52-card deck is shuffled and one card is drawn.
\begin{parts}
    \part[1] What is the probability of drawing a king?
    \part[1] What is the probability of drawing a red card?
    \part[2] What is the probability of drawing a king \textbf{or} a red card?
    \part[2] A card is drawn and \textbf{not} replaced; a second card is drawn. What is the
    probability both cards are aces?
\end{parts}
\begin{solution}[8cm]
\textbf{(a)} $P(\text{King}) = \dfrac{4}{52} = \dfrac{1}{13}$

\textbf{(b)} $P(\text{Red}) = \dfrac{26}{52} = \dfrac{1}{2}$

\textbf{(c)} $P(\text{King} \cup \text{Red}) = P(\text{King}) + P(\text{Red}) - P(\text{Red King})$
\[
= \frac{4}{52} + \frac{26}{52} - \frac{2}{52} = \frac{28}{52} = \mathbf{\frac{7}{13}}
\]

\textbf{(d)} Without replacement:
\[
P = \frac{4}{52} \times \frac{3}{51} = \frac{12}{2652} = \mathbf{\frac{1}{221}}
\]
\end{solution}

\question The table below shows the results of a survey of 200 students:

\begin{center}
\begin{tabular}{lccc}
\toprule
 & \textbf{Studies daily} & \textbf{Does not study daily} & \textbf{Total} \\
\midrule
\textbf{Passed} & 90 & 30 & 120 \\
\textbf{Failed} & 20 & 60 & 80 \\
\midrule
\textbf{Total} & 110 & 90 & 200 \\
\bottomrule
\end{tabular}
\end{center}

\begin{parts}
    \part[1] What is the probability a randomly selected student passed?
    \part[2] What is the probability a student studies daily \textbf{given} they passed?
    \part[2] What is the probability a student passed \textbf{given} they do not study daily?
    \part[3] Are the events ``passed'' and ``studies daily'' independent? Justify using
    probabilities.
\end{parts}
\begin{solution}[9cm]
\textbf{(a)} $P(\text{Pass}) = \dfrac{120}{200} = \mathbf{0.6}$

\textbf{(b)} $P(\text{Studies} \mid \text{Pass}) = \dfrac{90}{120} = \mathbf{0.75}$

\textbf{(c)} $P(\text{Pass} \mid \text{No Study}) = \dfrac{30}{90} = \mathbf{0.333}$

\textbf{(d)} Check: $P(\text{Pass}) \times P(\text{Studies}) = 0.6 \times \dfrac{110}{200}
= 0.6 \times 0.55 = 0.33$.\\
$P(\text{Pass} \cap \text{Studies}) = \dfrac{90}{200} = 0.45 \neq 0.33$.\\
Since $P(A \cap B) \neq P(A) \cdot P(B)$, the events are \textbf{not independent}.
\end{solution}

\question A box contains 4 defective and 6 good light bulbs. Two bulbs are drawn one at a time
\textbf{without replacement}.
\begin{parts}
    \part[3] Draw a tree diagram showing all possible outcomes and their probabilities.
    \part[2] Find the probability that exactly one bulb is defective.
    \part[2] Find the probability that at least one bulb is good.
\end{parts}
\begin{solution}[10cm]
\textbf{(a)} Let D = defective, G = good.
\[
P(DD) = \tfrac{4}{10}\cdot\tfrac{3}{9} = \tfrac{12}{90};\quad
P(DG) = \tfrac{4}{10}\cdot\tfrac{6}{9} = \tfrac{24}{90};\quad
P(GD) = \tfrac{6}{10}\cdot\tfrac{4}{9} = \tfrac{24}{90};\quad
P(GG) = \tfrac{6}{10}\cdot\tfrac{5}{9} = \tfrac{30}{90}
\]
(Total $= 90/90 = 1$\checkmark)

\textbf{(b)} $P(\text{exactly 1 defective}) = P(DG) + P(GD)
= \dfrac{24}{90} + \dfrac{24}{90} = \dfrac{48}{90} = \mathbf{\dfrac{8}{15}}$

\textbf{(c)} $P(\text{at least 1 good}) = 1 - P(DD) = 1 - \dfrac{12}{90}
= \dfrac{78}{90} = \mathbf{\dfrac{13}{15}}$
\end{solution}

\question In a certain town, 1\% of the population has a rare disease. A diagnostic test is
99\% accurate: it correctly identifies a positive case 99\% of the time, and correctly
identifies a negative case 99\% of the time.
\begin{parts}
    \part[2] Define events $D$ (has disease) and $T^+$ (tests positive). Write
    all given probabilities.
    \part[4] Using Bayes' Theorem, find the probability that a person \textbf{actually has
    the disease} given they \textbf{tested positive}. Show full working.
    \part[2] Comment on whether the test is reliable.
\end{parts}
\begin{solution}[10cm]
\textbf{(a)} $P(D) = 0.01$;\ $P(D^c) = 0.99$;\ $P(T^+\mid D) = 0.99$;\
$P(T^+\mid D^c) = 0.01$ (false positive rate).

\textbf{(b)} By the law of total probability:
\[
P(T^+) = P(T^+\mid D)P(D) + P(T^+\mid D^c)P(D^c)
= (0.99)(0.01) + (0.01)(0.99) = 0.0099 + 0.0099 = 0.0198
\]
By Bayes' Theorem:
\[
P(D \mid T^+) = \frac{P(T^+\mid D)\,P(D)}{P(T^+)}
= \frac{0.0099}{0.0198} = \mathbf{0.5}
\]

\textbf{(c)} Despite being 99\% accurate, a positive test result only means a 50\% chance of
actually having the disease. This is because the disease is so rare (base rate of 1\%) that
false positives are equally common as true positives.
\end{solution}

\end{questions}
\end{document}
